%%%%%%%%%%%%%%%%%%%%%%%%%%%%%%%%%%%%%%%%%%%%%%%%%
\section{Introdução}
%%%%%%%%%%%%%%%%%%%%%%%%%%%%%%%%%%%%%%%%%%%%%%%%%

A digitalização acelerada da indústria manufatureira tem gerado volumes crescentes de dados
provenientes de sensores distribuídos em equipamentos industriais. Esses dados, quando
devidamente coletados, armazenados e analisados, permitem detectar falhas precocemente,
otimizar a manutenção preditiva e aumentar a disponibilidade dos equipamentos
\cite{lee2015cyber}.

O presente trabalho descreve o desenvolvimento de uma plataforma completa de integração
de dados IoT aplicada ao contexto de monitoramento industrial. A solução integra dados
de múltiplas fontes heterogêneas — sensores em tempo real, sistemas ERP, APIs REST e
cadastros NoSQL — e os processa em um \textit{data lakehouse} estruturado em três camadas:
Bronze (dados brutos), Silver (dados limpos) e Gold (métricas agregadas)
\cite{armbrust2021lakehouse}.

A motivação central é demonstrar, na prática, os conceitos estudados na disciplina de
Integração de Dados II: modelagem NoSQL, arquitetura lakehouse, pipelines batch e
streaming, orquestração e governança de dados. O resultado é um sistema capaz de
processar mais de 100 eventos por segundo provenientes de sensores distribuídos em
três fábricas industriais, com detecção de anomalias em janelas de cinco minutos.

%%%%%%%%%%%%%%%%%%%%%%%%%%%%%%%%%%%%%%%%%%%%%%%%%
\section{Fundamentação Teórica}
%%%%%%%%%%%%%%%%%%%%%%%%%%%%%%%%%%%%%%%%%%%%%%%%%

\subsection{Arquitetura Medallion Lakehouse}

A arquitetura \textit{Medallion} organiza os dados em camadas progressivas de qualidade,
denominadas Bronze, Silver e Gold \cite{databricks2021medallion}. A camada Bronze armazena
os dados brutos exatamente como chegaram da fonte, sem transformações, garantindo
auditabilidade e a possibilidade de reprocessamento. A camada Silver aplica limpeza,
normalização e deduplicação, produzindo dados confiáveis para análise. A camada Gold
contém agregações e métricas de negócio prontas para consumo por ferramentas de
visualização e modelos de machine learning.

O conceito de \textit{lakehouse} unifica as vantagens dos \textit{data lakes} (baixo custo,
escalabilidade, suporte a dados semi-estruturados) com as dos \textit{data warehouses}
(suporte transacional, consultas SQL performáticas) \cite{armbrust2021lakehouse}.

\subsection{Processamento de Streams com Apache Kafka e Spark}

O Apache Kafka é uma plataforma distribuída de streaming de eventos que permite publicar,
assinar, armazenar e processar fluxos de registros em tempo real \cite{kreps2011kafka}.
Sua arquitetura de log imutável e particionamento por chave garante ordenação dentro de
partições e suporte a reprocessamento.

O Apache Spark Structured Streaming estende o modelo batch do Spark para processamento
contínuo, tratando streams como tabelas sem limites (\textit{unbounded tables})
\cite{zaharia2016apache}. O uso de \textit{watermarks} permite lidar com dados com
atraso (\textit{late data}), descartando eventos fora da janela de tolerância definida.

\subsection{Orquestração com Apache Airflow}

O Apache Airflow é uma plataforma de orquestração de workflows que representa pipelines
como grafos acíclicos direcionados (DAGs) \cite{airflow2024docs}. Cada nó do grafo é uma
tarefa com operadores especializados (PythonOperator, BashOperator, SparkSubmitOperator),
e as dependências entre tarefas definem a ordem de execução. O Airflow provê agendamento,
monitoramento, reexecução automática em caso de falha e registro histórico de execuções.

%%%%%%%%%%%%%%%%%%%%%%%%%%%%%%%%%%%%%%%%%%%%%%%%%
\section{Arquitetura da Plataforma}
%%%%%%%%%%%%%%%%%%%%%%%%%%%%%%%%%%%%%%%%%%%%%%%%%

A plataforma é composta por cinco camadas funcionais, todas containerizadas com Docker
Compose, conforme ilustrado na Figura~\ref{fig:arquitetura}.

\begin{figure}[h]
\centering
\begin{verbatim}
Sensores (Simulador Python)
        |
        v
   Apache Kafka ─────────────────────────────────┐
   tópico: iot-sensors                           | streaming
        |                                        v
        | consumer                    anomaly_detector.py
        v                             (janelas 5 min)
   MinIO / Bronze                            |
   factory_id/year/month/day                 v
        |                             MinIO / alerts/
        | Airflow @daily
        v
   Spark: bronze_to_silver.py
        |
        v
   MinIO / Silver (sensor_data/)
        |
        v
   Spark: silver_to_gold.py
        |
        v
   MinIO / Gold (Parquet)
        |
        v
   PostgreSQL (gold_equipment_metrics)
        |
        v
   Metabase (dashboards)

PostgreSQL (ERP) ── Airflow @daily  ──> Bronze/erp/
API Manutenção   ── Airflow @hourly ──> Bronze/maintenance/
MongoDB          ── cadastro de equipamentos e sensores
\end{verbatim}
\caption{Arquitetura da plataforma IoT Industrial}
\label{fig:arquitetura}
\end{figure}

\subsection{Camada de Fontes de Dados}

Quatro fontes alimentam a plataforma:

\begin{itemize}
    \item \textbf{Sensores IoT}: simulador Python gera eventos JSON em tempo real para o
    tópico Kafka \texttt{iot-sensors}, com 100 eventos/segundo e taxa de anomalia
    configurável (padrão 5\%). Cada evento inclui tipo de medição (temperatura, umidade,
    pressão, vibração, corrente), valor, unidade, qualidade e metadados do sensor.

    \item \textbf{Sistema ERP (PostgreSQL)}: contém registros de manutenção e ordens
    de trabalho dos equipamentos. A ingestão batch diária é orquestrada pelo Airflow.

    \item \textbf{API de Manutenção (REST)}: fornece dados de ordens de manutenção
    agendadas. A ingestão horária é realizada via PythonOperator no Airflow, com
    fallback para geração de dados \textit{mock} caso a API esteja indisponível.

    \item \textbf{Cadastro de Equipamentos (MongoDB)}: armazena 50 equipamentos
    distribuídos em 3 fábricas (SP, RJ, MG), com sensores e ranges normais embutidos
    no documento do equipamento, conforme decisão arquitetural ADR-001.
\end{itemize}

\subsection{Camada de Armazenamento — Medallion Lakehouse}

Todo o armazenamento de dados utiliza o MinIO, sistema de armazenamento de objetos
compatível com S3, estruturado em cinco buckets:

\begin{itemize}
    \item \textbf{bronze}: dados brutos em formato JSONL, particionados por
    \texttt{factory\_id/year/month/day}. Acumula lotes de 1.000 eventos por fábrica,
    além de subpastas \texttt{erp/} e \texttt{maintenance/} para dados de outras fontes.

    \item \textbf{silver}: dados de sensores limpos, normalizados e deduplicados,
    em formato JSONL no caminho \texttt{sensor\_data/}.

    \item \textbf{gold}: métricas horárias por equipamento em formato Parquet,
    particionadas por \texttt{year/month/factory\_id}, no caminho
    \texttt{equipment\_metrics\_hourly/}.

    \item \textbf{alerts}: alertas de anomalias gerados pelo pipeline de streaming,
    particionados por \texttt{factory\_id}.

    \item \textbf{checkpoints}: checkpoints do Spark Structured Streaming para
    garantia de processamento exatamente-uma-vez (\textit{exactly-once}).
\end{itemize}

\subsection{Camada de Processamento}

\textbf{Bronze → Silver} (\texttt{bronze\_to\_silver.py}): job Spark que lê todos os
arquivos JSONL do bucket Bronze com schema explícito do sensor, filtra apenas eventos
com \texttt{sensor\_id} não nulo (descartando dados de ERP e manutenção), aplica
limpeza (trim, lowercase, conversão de timestamp), extrai campos de metadados e
remove duplicatas por \texttt{event\_id}.

\textbf{Silver → Gold} (\texttt{silver\_to\_gold.py}): job Spark que lê os dados da
camada Silver e produz agregações horárias por equipamento e tipo de medição, calculando
média, mínimo, máximo, desvio padrão, contagem de eventos, taxa de anomalia e
\textit{criticality score} (combinação ponderada de taxa de anomalia e variabilidade).

\textbf{Detecção de Anomalias} (\texttt{anomaly\_detector.py}): job Spark Structured
Streaming que consome o tópico Kafka em tempo real, aplica janelas tumbantes de cinco
minutos com \textit{watermark} de dez minutos para tolerância a dados atrasados, e
classifica alertas em dois níveis: \textit{warning} (taxa de anomalia ≥ 10\%) e
\textit{critical} (taxa ≥ 30\%).

\subsection{Camada de Orquestração}

O Apache Airflow 2.9.3 com LocalExecutor gerencia três DAGs:

\begin{itemize}
    \item \textbf{iot\_lakehouse\_pipeline} (@daily): executa sequencialmente
    Bronze→Silver via BashOperator com \texttt{spark-submit}, seguido de Silver→Gold.
    Configurado com 2 retentativas e intervalo de 5 minutos entre tentativas.

    \item \textbf{erp\_batch\_ingestion} (@daily): extrai registros de eventos de
    equipamentos do PostgreSQL via PythonOperator e grava em \texttt{bronze/erp/}.

    \item \textbf{maintenance\_api\_ingestion} (@hourly): consome a API REST de
    manutenção e grava em \texttt{bronze/maintenance/}.
\end{itemize}

%%%%%%%%%%%%%%%%%%%%%%%%%%%%%%%%%%%%%%%%%%%%%%%%%
\section{Implementação}
%%%%%%%%%%%%%%%%%%%%%%%%%%%%%%%%%%%%%%%%%%%%%%%%%

\subsection{Stack Tecnológica}

A Tabela~\ref{tab:stack} apresenta todas as tecnologias utilizadas e suas versões.

\begin{table}[h]
\centering
\caption{Stack tecnológica da plataforma}
\label{tab:stack}
\begin{tabular}{|l|l|l|}
\hline
\textbf{Categoria} & \textbf{Tecnologia} & \textbf{Versão} \\
\hline
Containerização & Docker + Docker Compose & 24.x / 2.x \\
\hline
Streaming & Apache Kafka (Confluent) & 7.5.0 \\
\hline
Processamento & Apache Spark & 3.5.3 \\
\hline
Orquestração & Apache Airflow & 2.9.3 \\
\hline
Object Storage & MinIO & latest \\
\hline
NoSQL & MongoDB & 6.x \\
\hline
SQL & PostgreSQL & 15 \\
\hline
Visualização & Metabase & 0.49.0 \\
\hline
Linguagem & Python & 3.10+ \\
\hline
\end{tabular}
\end{table}

\subsection{Modelagem NoSQL — MongoDB}

A decisão de modelagem mais relevante foi a desnormalização do cadastro de equipamentos:
sensores e seus ranges de medição são embutidos diretamente no documento do equipamento,
conforme ADR-001. Essa decisão elimina \textit{joins} no pipeline de streaming, que precisa
consultar os ranges normais para detectar anomalias em tempo real. O custo é uma
duplicação caso o mesmo sensor esteja em múltiplos equipamentos — cenário inexistente
no domínio do projeto.

Os índices criados incluem \texttt{factory\_id}, \texttt{type}, \texttt{status},
\texttt{sensors.id}, \texttt{sensors.type} (em arrays) e o índice composto
\texttt{\{factory\_id, status\}} para a consulta mais frequente no pipeline.

\subsection{Particionamento dos Dados}

A estratégia de particionamento foi otimizada para as queries downstream:

\textbf{Bronze}: \texttt{factory\_id/year/month/day} — \texttt{factory\_id} como
primeira dimensão permite que o Spark aplique \textit{partition pruning} ao processar
uma fábrica específica, reduzindo I/O em aproximadamente 67\% para queries de fábrica
única. Lotes de 1.000 eventos por fábrica evitam o problema de \textit{small files}.

\textbf{Silver}: escrita flat em \texttt{sensor\_data/} com \texttt{coalesce(4)},
otimizando para leitura completa pelo job Silver→Gold.

\textbf{Gold}: \texttt{year/month/factory\_id} em formato Parquet, permitindo
\textit{partition pruning} por fábrica e mês nas queries analíticas do Metabase.

\subsection{Qualidade de Dados}

Dez \textit{assertions} automatizadas foram implementadas com pytest, cobrindo:

\begin{enumerate}
    \item Presença de campos obrigatórios (\texttt{event\_id}, \texttt{sensor\_id}, etc.)
    \item Validade do tipo de medição (enum com 5 valores)
    \item Validade do campo \texttt{quality} (good/warning/bad)
    \item \texttt{factory\_id} pertencente ao conjunto conhecido de fábricas
    \item Valores dentro dos ranges físicos por tipo de sensor
    \item Taxa de duplicatas abaixo de 1\% (Bronze é append-only)
    \item \texttt{is\_anomaly} do tipo booleano
    \item \texttt{value} do tipo numérico
    \item \texttt{timestamp} não vazio
    \item \texttt{event\_id} no formato UUID v4
\end{enumerate}

%%%%%%%%%%%%%%%%%%%%%%%%%%%%%%%%%%%%%%%%%%%%%%%%%
\section{Guia Operacional}
%%%%%%%%%%%%%%%%%%%%%%%%%%%%%%%%%%%%%%%%%%%%%%%%%

\subsection{Pré-requisitos}

\begin{itemize}
    \item Docker Desktop 24+ e Docker Compose 2.x instalados e em execução
    \item Python 3.10+
    \item Dependências: \texttt{pip install kafka-python pymongo boto3 psycopg2-binary requests pytest}
\end{itemize}

\subsection{Subindo o Ambiente}

\begin{verbatim}
# 1. Subir todos os containers
docker-compose up -d

# 2. Aguardar containers ficarem healthy (~3-5 minutos)
docker-compose ps

# 3. Criar banco de dados Airflow e Metabase (primeira execução)
docker exec postgres psql -U user -d iot -c "CREATE DATABASE airflow;"
docker exec postgres psql -U user -d iot -c "CREATE DATABASE metabase;"

# 4. Popular MongoDB
python src/ingestao/mongodb_setup.py

# 5. Criar buckets no MinIO
docker exec minio mc alias set local http://localhost:9000 admin password
docker exec minio mc mb local/bronze local/silver local/gold \
    local/alerts local/checkpoints
\end{verbatim}

\subsection{Credenciais de Acesso}

A Tabela~\ref{tab:creds} lista as credenciais de acesso a todos os serviços da plataforma.

\begin{table}[h]
\centering
\caption{Credenciais de acesso aos serviços}
\label{tab:creds}
\begin{tabular}{|l|l|l|l|}
\hline
\textbf{Serviço} & \textbf{URL} & \textbf{Usuário} & \textbf{Senha} \\
\hline
MinIO Console & http://localhost:9001 & admin & password \\
\hline
Apache Airflow & http://localhost:8080 & admin & admin \\
\hline
Metabase & http://localhost:3030 & (primeiro acesso) & (primeiro acesso) \\
\hline
Spark Master UI & http://localhost:8081 & — & — \\
\hline
PostgreSQL & localhost:5432 & user & password \\
\hline
MongoDB & localhost:27017 & — & — \\
\hline
Kafka & localhost:9092 & — & — \\
\hline
\end{tabular}
\end{table}

\subsection{Executando o Pipeline Completo}

\begin{verbatim}
# Terminal 1 — Geração de dados de sensores
python src/ingestao/sensor_simulator.py --duration 120

# Terminal 2 — Consumer Kafka → Bronze
python src/streaming/kafka_to_bronze.py

# Após gerar dados suficientes, executar via Airflow:
# http://localhost:8080 → iot_lakehouse_pipeline → Trigger DAG

# Ou executar manualmente:
docker exec spark-master /opt/spark/bin/spark-submit \
  --master spark://spark-master:7077 \
  --conf spark.jars.ivy=/tmp/ivy2 \
  --packages org.apache.hadoop:hadoop-aws:3.3.4,\
com.amazonaws:aws-java-sdk-bundle:1.12.262 \
  /opt/spark/jobs/bronze_to_silver.py

docker exec spark-master /opt/spark/bin/spark-submit \
  --master spark://spark-master:7077 \
  --conf spark.jars.ivy=/tmp/ivy2 \
  --packages org.apache.hadoop:hadoop-aws:3.3.4,\
com.amazonaws:aws-java-sdk-bundle:1.12.262 \
  /opt/spark/jobs/silver_to_gold.py
\end{verbatim}

\subsection{Detecção de Anomalias em Streaming}

\begin{verbatim}
docker exec spark-master /opt/spark/bin/spark-submit \
  --master spark://spark-master:7077 \
  --conf spark.jars.ivy=/tmp/ivy2 \
  --packages org.apache.spark:spark-sql-kafka-0-10_2.12:3.5.0,\
org.apache.hadoop:hadoop-aws:3.3.4,\
com.amazonaws:aws-java-sdk-bundle:1.12.262 \
  /opt/spark/jobs/anomaly_detector.py
\end{verbatim}

\subsection{Testes de Qualidade}

\begin{verbatim}
pip install pytest
pytest tests/ -v
# Esperado: 10 testes passando
\end{verbatim}

\subsection{Exportando Gold para o Metabase}

\begin{verbatim}
docker exec spark-master /opt/spark/bin/spark-submit \
  --master spark://spark-master:7077 \
  --conf spark.jars.ivy=/tmp/ivy2 \
  --packages org.apache.hadoop:hadoop-aws:3.3.4,\
com.amazonaws:aws-java-sdk-bundle:1.12.262,\
org.postgresql:postgresql:42.6.0 \
  /opt/spark/jobs/gold_to_postgres.py
\end{verbatim}

Após a exportação, acessar o Metabase em \url{http://localhost:3030}, sincronizar
o schema do banco \texttt{iot} e a tabela \texttt{gold\_equipment\_metrics} estará
disponível para criação de dashboards.

\subsection{Parando o Ambiente}

\begin{verbatim}
# Para mas mantém dados (volumes preservados)
docker-compose down

# Para e remove todos os dados
docker-compose down -v
\end{verbatim}

%%%%%%%%%%%%%%%%%%%%%%%%%%%%%%%%%%%%%%%%%%%%%%%%%
\section{Resultados}
%%%%%%%%%%%%%%%%%%%%%%%%%%%%%%%%%%%%%%%%%%%%%%%%%

A plataforma foi implementada com sucesso, atendendo a todos os requisitos das
Entregas 1, 2 e 3 do projeto integrador. Os principais resultados obtidos foram:

\begin{itemize}
    \item \textbf{Ingestão}: o simulador de sensores gerou e publicou eventos no Kafka
    a uma taxa média de 100 eventos/segundo, distribuídos entre 50 equipamentos de
    3 fábricas. O consumer Kafka gravou lotes de 1.000 eventos por fábrica no bucket
    Bronze do MinIO, totalizando aproximadamente 180.000 eventos em uma execução de
    30 minutos.

    \item \textbf{Processamento Batch}: os jobs Spark Bronze→Silver e Silver→Gold
    executaram com sucesso via Airflow e de forma manual, processando o volume de dados
    acumulado em menos de 2 minutos cada. A camada Gold produz métricas horárias
    com \textit{criticality score} para cada combinação de equipamento e tipo de medição.

    \item \textbf{Múltiplas Fontes}: além dos sensores, a plataforma integrou com
    sucesso dados do ERP PostgreSQL (190 registros de eventos de equipamentos) e
    da API REST de manutenção (dados \textit{mock} com 10--30 registros/hora).

    \item \textbf{Qualidade}: todos os 10 testes automatizados passaram na camada
    Bronze, validando schema, tipos, ranges físicos e integridade referencial dos dados.

    \item \textbf{Visualização}: dados do Gold foram exportados para o PostgreSQL
    e visualizados no Metabase com dashboards de métricas por fábrica, tipo de evento
    e evolução temporal.
\end{itemize}

%%%%%%%%%%%%%%%%%%%%%%%%%%%%%%%%%%%%%%%%%%%%%%%%%
\section{Conclusão}
%%%%%%%%%%%%%%%%%%%%%%%%%%%%%%%%%%%%%%%%%%%%%%%%%

Este trabalho demonstrou a viabilidade de construir uma plataforma completa de
integração de dados IoT industrial utilizando ferramentas open-source amplamente
adotadas pela indústria. A arquitetura Medallion Lakehouse provou-se adequada para
o domínio, separando claramente responsabilidades entre camadas e permitindo reprocessamento
de dados históricos sem perda dos dados brutos.

Os principais desafios enfrentados foram a integração do Spark com o MinIO via protocolo
S3A dentro de containers Docker (exigindo configurações específicas de endpoint e
\textit{path style access}), e a gestão de múltiplos schemas heterogêneos no mesmo
bucket Bronze — resolvida com leitura por schema explícito e filtragem por \texttt{sensor\_id}.

Como trabalhos futuros, destacam-se: (i) adoção de formatos de tabela como Apache Iceberg
para suporte transacional ACID no lakehouse; (ii) implementação de CDC (\textit{Change
Data Capture}) no PostgreSQL ERP para ingestão incremental; e (iii) desenvolvimento de
modelos preditivos de manutenção utilizando as métricas da camada Gold.

%%%%%%%%%%%%%%%%%%%%%%%%%%%%%%%%%%%%%%%%%%%%%%%%%
\section*{Agradecimentos}
%%%%%%%%%%%%%%%%%%%%%%%%%%%%%%%%%%%%%%%%%%%%%%%%%

Os autores agradecem ao corpo docente do curso de Ciência de Dados e Inteligência
Artificial da PUC Minas pelo suporte e orientação durante o desenvolvimento deste
projeto.
